% LSPV-1 REPORT::B3-DICHOTOMY-ER-CONTAINMENT345::2026-11-02T16:00-08:00
\documentclass[11pt]{article}
\usepackage[margin=0.95in]{geometry}
\usepackage{amsmath,amssymb,amsthm,hyperref}
\newtheorem{theorem}{Theorem}
\newtheorem{proposition}[theorem]{Proposition}
\newtheorem{corollary}[theorem]{Corollary}
\title{B3-DICHOTOMY-ER-CONTAINMENT345:\\ What the Worst-Case Derivation Bridge Can and Cannot Be}
\author{PNP General SAT Lower-Bound Lead (LSPV-1)}
\date{November 2, 2026, 16:00 PST}
\begin{document}\maketitle
\noindent\textbf{Status.} PROVED/EXACT: Theorems 1--3 (elementary, in the declared measure). KNOWN THEOREM: Haken 1985; Cook--Reckhow 1979; Cook 1975 / Krajicek--Pudlak translation of PV proofs into extended Frege; Tseitin extension. CANDIDATE: $R_{\rm mix341}\equiv_p eF$. OPEN: B3 and a superpolynomial $\overline W_R$ for $R_{\rm mix341}$. \textbf{P versus NP remains OPEN.} No external referee has checked this report.

\section{Fixed objects}
$\overline W_R(N)=\max_{F\in{\rm UNSAT},|\langle F\rangle|\le N}\min_{\pi:F\vdash_R\bot}{\rm cost}_R(\pi)$, with the optimal $\pi$ granted free. Write
\[
\mathrm{L}(R):\ \overline W_R\ \text{is superpolynomial},\qquad
\mathrm{B3}(R):\ \mathrm{SAT}\in\mathrm P\Rightarrow\overline W_R\le N^{O(1)}.
\]
Then $\mathrm{L}(R)\wedge\mathrm{B3}(R)\Rightarrow \mathrm{SAT}\notin\mathrm P$ (PROVED/EXACT, pure logic). Cost is proof size (number of charged nodes/clauses); for two systems the comparisons below use cost-preserving polynomial simulations.

\section{Three exact theorems}
\begin{theorem}[Monotonicity]\label{t:mono}
If $R'$ polynomially simulates $R$ in the cost measure, then $\mathrm{L}(R')\Rightarrow\mathrm{L}(R)$ and $\mathrm{B3}(R)\Rightarrow\mathrm{B3}(R')$. Hence lower bounds are easiest for weak systems, and the bridge is easiest for strong systems; the two arrows pull in opposite directions.
\end{theorem}
\begin{proof} $\overline W_{R'}\le{\rm poly}(\overline W_R)$ by the simulation. A superpolynomial $\overline W_{R'}$ forces a superpolynomial $\overline W_R$; a polynomial $\overline W_R$ forces a polynomial $\overline W_{R'}$.\end{proof}

\begin{theorem}[Dichotomy: the bridge is vacuous or open]\label{t:dich}
If $\mathrm L(R)$ is a theorem, then $\mathrm{B3}(R)\iff \mathrm{SAT}\notin\mathrm P$, i.e.\ $\mathrm P\ne\mathrm{NP}$. This applies to ordinary Resolution (Haken 1985), constant-depth Frege, cutting planes and every other system with an unconditional superpolynomial bound.
\end{theorem}
\begin{proof} $\mathrm{B3}(R)$ is ``$A\Rightarrow C$'' with $C$ false, so it is equivalent to $\neg A$. \end{proof}
\noindent So the ``weakest sufficient bridge'' is informative only for a system $R$ whose lower bound is \emph{not} yet known. For such $R$ the pair $(\mathrm L,\mathrm{B3})$ is jointly sufficient but not necessary for $\mathrm P\ne\mathrm{NP}$: a lower bound for one fixed $R$ does not even imply $\mathrm{NP}\ne\mathrm{coNP}$ (Cook--Reckhow need every system), and $\mathrm{NP}=\mathrm{coNP}$ is consistent with $\mathrm P\ne\mathrm{NP}$ for all that is known.

\begin{theorem}[Containment of extended resolution]\label{t:er}
Let $R_{\rm mix341}$ be the calculus of LSPV-2 report 341: unrestricted-width binary RES at unit cost, constant-arity AND3, affine XOR and DEF, where DEF introduces one fresh variable $z\leftrightarrow g(x,y)$ for a two-input Boolean gate at unit cost. Then $R_{\rm mix341}$ polynomially simulates Extended Resolution (ER) in the node-count measure, so $\overline W_{R_{\rm mix341}}\le \overline W_{\rm ER}\cdot O(1)$, and
\[
\mathrm L(R_{\rm mix341})\ \Longrightarrow\ \mathrm L(\mathrm{ER})\iff \mathrm L(eF).
\]
\end{theorem}
\begin{proof} Every ER extension $z\leftrightarrow\varphi$ is replaced by a chain of binary-gate definitions of the subformulas of $\varphi$ (Tseitin), with the three defining clauses per gate charged to one DEF node; the number of nodes grows by the size of $\varphi$. ER and eF are polynomially equivalent (Cook--Reckhow 1979; Kullmann/Tseitin), which gives the final equivalence. \end{proof}

\begin{corollary}[Honest status of the target]
A superpolynomial worst-case lower bound for the project calculus $R_{\rm mix341}$ is at least as strong as a superpolynomial extended Frege lower bound, a problem open since 1979 and widely regarded as beyond current techniques (Pich--Santhanam, ``Why are proof complexity lower bounds hard?'', FOCS 2019). The extra rules AND3 and affine XOR are sound and locally checkable; whether eF simulates them (so that $R_{\rm mix341}\equiv_p eF$) is a CANDIDATE that needs a per-rule audit, in particular the allowed width of an affine combination.
\end{corollary}

\section{What a real bridge would look like}
\textbf{KNOWN THEOREM (Cook 1975; Krajicek--Pudlak).} If $A$ is a polynomial-time SAT decider and PV proves the Pi-1 statement ``$A(F)=0\Rightarrow\neg{\rm Sat}(F,a)$'' for all $F,a$, then its propositional translations have polynomial-size eF proofs, and substituting the polynomial-size computation of $A$ on a given unsatisfiable $F$ gives an eF refutation of size $q_A(N+T_A^{wc}(N))$. This is exactly the \emph{shape} of B2.

\begin{proposition}[Sufficient reflection condition]\label{p:ref}
Let $R\supseteq eF$ (e.g.\ $R_{\rm mix341}$ if Theorem~\ref{t:er} and the CANDIDATE hold). If for every correct polynomial-time SAT decider $A$ the reflection family $\mathrm{Ref}_A$ (``a valid accepting-UNSAT trace of $A$ on $F$ implies $\neg F$'') has polynomial-size $R$-proofs, then $\mathrm{B3}(R)$ holds, with $\overline W_R(N)\le q_A(N+T_A^{wc}(N))$ for each $A$.
\end{proposition}
\begin{proof} Instantiate $\mathrm{Ref}_A$ at $F$ and add the trace, whose size is $T_A(F)$ and which $R$ verifies by polynomially many local steps. \end{proof}
\noindent This is the ``correctness-reflection compiler'' that LSPV-4 345 and LSPV-1 338 name as the missing arrow. The hypothesis is not supplied by $\mathrm{SAT}\in\mathrm P$ alone: $A$ may be correct without $R$ (or PV) being able to prove it. Proposition~\ref{p:ref} is a sufficient condition only, not an equivalence. The unproved implication ``$\mathrm P=\mathrm{NP}\Rightarrow eF$ polynomially bounded'' stays open and is a famous problem, not a project-specific gap.

\section{Consequences for the five-agent plan}
\begin{enumerate}
\item \textbf{Do not spend effort on B3 for Resolution-like fragments.} By Theorem~\ref{t:dich} it is equivalent to the target.
\item \textbf{Hard-instance work (LSPV-2/3) should be scored against eF.} A family hard for $R_{\rm mix341}$ is a family hard for eF; any claimed family should be tested first against the known short eF proofs of pigeonhole, Tseitin and counting principles (Buss 1987; Cook--Reckhow).
\item \textbf{Strategy-search remains free but irrelevant.} Granting the optimal derivation cannot rescue a weak calculus: Theorem~\ref{t:dich} holds with free search.
\item \textbf{Restrict the target.} The only non-vacuous restricted bridge on record is the PV-provable-solver case above. A superpolynomial bound in that restricted class would be a different, conditional statement, and must be labeled CONDITIONAL.
\end{enumerate}

\section{Literature and audit}
\textbf{CORPUS REVIEWED:} AGENTS.md; LSPV-1 338 (moving-decider minimax, B2 reformulation); LSPV-2 341 (calculus $R_{\rm mix341}$, DEF rule at line 20); LSPV-4 344 (space-capped portfolio); MAIN REFEREE297; root \texttt{main.tex} LSPV-1 region (lines 679--3570). Search this run: proof-complexity literature on whether $\mathrm P=\mathrm{NP}$ implies a polynomially bounded eF (Cook--Reckhow equivalence $\mathrm{NP}=\mathrm{coNP}\iff$ some polybounded system; eF bounded-ness open) and Pich--Santhanam on conditional barriers; recent hits: Khaniki--Pich--Sokolov (CCC 2026) and ``The Proof Analysis Problem'' (FOCS 2025), both about the difficulty of analyzing eF lower bounds. Searches were standard-mode web results only; the PV-translation statement is cited from standard references, not re-derived from a source read this run.

\textbf{DUPLICATION / CONTRADICTION AUDIT:} Monotonicity and Theorem~\ref{t:dich} are elementary and probably folklore; no novelty claim. The ER-containment is stated from the rules in 341; it contradicts no earlier claim but sharpens the 341 remark that the linear bound $q\ge m-1$ is the only proved lower bound for $R_{\rm mix341}$.

\textbf{WHAT CHANGED RELATIVE TO ROOT MAIN:} adds the dichotomy and the eF-containment status of the target; changes no result.

\textbf{ADOPTED / REJECTED / RETIRED:} adopted B2/B3 definitions of 338; rejected B3 for known-lower-bound systems as a stepping stone (equivalent to the goal); retired any reading of the shortest-DAG framework as weaker than an eF lower bound.

\textbf{STALE OR OVERSTATED CLAIMS FOUND:} descriptions calling B3 the ``weakest sufficient bridge'' are correct only for systems with unknown lower bounds; for others the bridge equals $\mathrm P\ne\mathrm{NP}$.

\textbf{Time-space audit (LSPV-4 role).} No computation was run. All statements are symbolic; proof-size accounting is by node count, with no space measure claimed.

\textbf{Next exact targets.} (i) Per-rule audit of AND3 and affine XOR for eF-simulation (decides $R_{\rm mix341}\equiv_p eF$). (ii) Decide whether a nonrelativizing argument can derive the reflection family $\mathrm{Ref}_A$ in $R$ from $\mathrm{SAT}\in\mathrm P$ alone.

\begin{thebibliography}{9}
\bibitem{cr} S. Cook, R. Reckhow, The relative efficiency of propositional proof systems, JSL 44 (1979).
\bibitem{h} A. Haken, The intractability of resolution, TCS 39 (1985).
\bibitem{ps} O. Pich, R. Santhanam, Why are proof complexity lower bounds hard?, FOCS 2019.
\bibitem{k} J. Krajicek, Bounded Arithmetic, Propositional Logic and Complexity Theory, CUP 1995.
\end{thebibliography}
\end{document}
% END_OF_VERIFIED_REPORT::LSPV-1::B3-DICHOTOMY-ER-CONTAINMENT345
